\section{Probability limits and squarefree restrictions}\label{sec:probability}

\subsection{The unsigned law in each progression}

\begin{proposition}\label{prop:unsigned}
Under Assumption~\ref{ass:inputs}, $V(p)$ converges in each fixed
finite set of coordinates to $\PD(1)$ when $p$ ranges over a fixed
reduced progression, in either the dyadic or cumulative convention.
\end{proposition}

\begin{proof}
First take $p$ uniform in $(x,2x]$ in the class. The denominator
$N_M(x)$ is asymptotic to $x/(\varphi(M)L)$ by the fixed-modulus
prime number theorem. For rectangles $\prod_{i=1}^d(a_i,b_i]$ with
positive lower endpoints and $\sum b_i<1$, choose prime slots
$(x^{a_i},x^{b_i}]$. Their reciprocal tuple sum tends to
$\prod_i\log(b_i/a_i)$; coincidences contribute $o(1)$.
Equation~\eqref{eq:interiorunsigned} and smooth approximation of
the dyadic cutoff give the same limit for the mean tuple count.

Regard the prime factors of $p-1$, with multiplicity, as labelled
balls of masses $\log q/\log(p-1)$. On a compact subset of
\[
 \mathcal D_d=\{(t_1,\ldots,t_d):t_i>0,\ \sum_i t_i<1\},
\]
the expected ordered distinct-ball measure therefore converges to
$\prod_i\dd t_i/t_i$. The change from $\log x$ to $\log(p-1)$
is $1+O(1/L)$ and is controlled by endpoint strips. Repeated
numerical primes on a fixed interior compact set entail
$q^2\mid p-1$ with $q\geq x^{a/2}$ for some fixed $a>0$.
Their integer count is
$O(x^{1-a/2}+\sqrt{x})=o(N_M(x))$. Tuple counts are bounded on
such a compact set, so multiplicities cause no further error.

Sample balls without replacement with probability proportional to
their remaining masses. The first $d$ size-biased draws multiply
the factorial density by
\[
 \prod_{i=1}^d\frac{t_i}{1-\sum_{j<i}t_j}.
\]
The resulting density is
$h_d(t)=\prod_{i=1}^d(1-\sum_{j<i}t_j)^{-1}$ on $\mathcal D_d$.
The triangular substitution
$u_i=t_i/(1-\sum_{j<i}t_j)$ identifies it with uniform measure
on $(0,1)^d$. Its mass is one. Compact exhaustion therefore excludes
escaped boundary mass and proves convergence of the size-biased
draws to the first $d$ $\GEM(1)$ fragments.

Let $R_{x,d}$ be the unseen mass after $d$ draws, with zero padding
after exhaustion. Its expectation tends to $2^{-d}$.
Sorting the first $d$ draws and padding by zero approximates every
ranked coordinate of $V(p)$ within $R_{x,d}$.
Uniform continuity and Markov's inequality, followed by $d\to\infty$,
prove the ranked finite-dimensional law.
A finite decomposition of $(X/2^J,X]$ into dyads gives the cumulative
law. The omitted initial class-prime proportion tends to $2^{-J}$;
take $X\to\infty$ and then $J\to\infty$.
\end{proof}

\subsection{Complex measures and Fourier inversion}

Attach $f_z(p-1)$ to each sampled prime, for a fixed $z\ne1$.
The resulting complex factorial measure $\sigma_{x,d}^{z}$ has
total variation dominated by its unsigned counterpart.
Equation~\eqref{eq:interiorsigned} implies
\[
 \int\phi\dd\sigma_{x,d}^{z}\longrightarrow0,
       \qquad\phi\in C_c(\mathcal D_d).
\]
The same unsigned endpoint-strip and multiplicity bounds used above
justify the normalization and rectangle-to-continuous-test passage
in absolute value.

Let $\nu_{x,d}$ be the unsigned size-biased law and
$\nu_{x,d}^{z}$ its complex weighted version. On the interior,
\[
 \dd\nu_{x,d}^{z}
  =\prod_i\frac{t_i}{1-\sum_{j<i}t_j}\dd\sigma_{x,d}^{z},
 \qquad |\nu_{x,d}^{z}|\leq\nu_{x,d}.
\]
Choose $0\leq\chi\leq1$ compactly supported in $\mathcal D_d$
with $\int\chi h_d>1-\eta$. For any continuous $H$ on the cube,
\[
 \left|\int H\dd\nu_{x,d}^{z}\right|
 \leq\left|\int\chi H\dd\nu_{x,d}^{z}\right|
       +\|H\|_\infty\int(1-\chi)\dd\nu_{x,d}.
\]
The first term tends to zero, and the limsup of the second is at
most $\eta\|H\|_\infty$. Let $\eta\to0$. Thus
$\nu_{x,d}^{z}$ converges weakly to the zero complex measure on
the whole cube. This is weak convergence, not convergence in total
variation. At $d=1$ its total mass is exactly $\E f_z(p-1)$,
which proves cancellation without a separate $d=0$ correlation.

If $H$ tests $k$ ranked coordinates, the error from sorting $d$
draws, with $d\geq k$, is at most
\[
 \omega_H(\eta)+2\|H\|_\infty\E R_{x,d}/\eta.
\]
The same bound holds for the complex weighted average because
$|f_z|=1$. Take $x\to\infty$, then $d\to\infty$, then
$\eta\to0$. We obtain
\begin{equation}\label{eq:modecancel}
 \E[f_z(p-1)H(V_1(p),\ldots,V_k(p))]\longrightarrow0
       \quad(z\ne1).
\end{equation}
For fixed $m$, finite Fourier inversion gives
\[
 \ind_{\Omega(n)\equiv a\ (m)}
  =\frac1m\sum_{j=0}^{m-1}
       \exp(-2\pi ija/m)\exp(2\pi ij\Omega(n)/m).
\]
The trivial mode is Proposition~\ref{prop:unsigned}; $z=-1$ is
the parity case; the other modes were treated in
Section~\ref{sec:model}. This proves Theorem~\ref{thm:main} in
dyads. The same finite dyadic decomposition proves its cumulative
form.

\subsection{Local square exclusions}

Let $B_Y(n)$ indicate that $q^2\nmid n$ for every prime $q\leq Y$.
Within a fixed reduced class modulo $M$, this condition on $p-1$
is a finite union of reduced classes modulo
$\lcm(M,\prod_{q\leq Y}q^2)$, apart from finitely many primes.
For fixed $Y$, apply the unsigned law and \eqref{eq:modecancel}
with $z=-1$ in each class. Reduced-lift counting gives the factors
$1-\rho_q(M,b)$ in \eqref{eq:rho}.
If $q\nmid M$, there are $q(q-1)$ reduced lifts modulo $q^2$
and one has $p\equiv1\pmod{q^2}$. If $v_q(M)=1$ and
$b\equiv1\pmod q$, one of $q$ lifts has that property. If
$q^2\mid M$, it is forced or excluded by $b$ itself.

To let $Y$ grow after $x$, use
\begin{equation}\label{eq:squaretail}
 \limsup_{x\to\infty}
 \frac{\#\{x<p\leq2x:\exists q>Y, q^2\mid p-1\}}
      {\pi(2x)-\pi(x)}\ll1/Y.
\end{equation}
For $Y<q\leq x^{1/4}$, Brun--Titchmarsh (in the form recorded in \cite{Ford}) gives
$O(x/[q(q-1)\log x])$ primes in the relevant progression.
For $q>x^{1/4}$, count integers and use
\[
 \sum_{x^{1/4}<q\leq\sqrt{2x}}(x/q^2+1)
       =O(x^{3/4}+\sqrt{x})=o(x/\log x).
\]
The first reciprocal tail is $O(1/Y)$.
Changing the denominator in \eqref{eq:squaretail} to $N_M(x)$
costs only a fixed factor. Bounded tests and signs preserve the bound.

Write $B=\mu^2$ temporarily for the arithmetic indicator.
Taking $x\to\infty$ and then $Y\to\infty$ yields the four limits
\begin{align*}
 \E H(V(p))&\longrightarrow\E H(L),&
 \E[\lambda(p-1)H(V(p))]&\longrightarrow0,\\
 \E[B(p-1)H(V(p))]&\longrightarrow C_{M,b}\E H(L),&
 \E[\lambda(p-1)B(p-1)H(V(p))]&\longrightarrow0.
\end{align*}
The identities
\[
 \ind_{\lambda=s,B=1}=B(1+s\lambda)/2,\qquad
 \ind_{\lambda=s,B=0}=(1-B)(1+s\lambda)/2
\]
give every joint atom. This proves Theorem~\ref{thm:squarefree},
including $C_{M,b}=0$. Finally $\mu=\lambda\mu^2$ gives
\eqref{eq:mu} and the two M\"obius-sign densities.

\subsection{The prescribed-valuation consequence}

On the event $v_q(p-1)=e_q$ for $q\in\mathcal S$, let
$B_{\mathcal S}$ indicate squarefreeness outside $\mathcal S$.
Then
\[
 \mu_{\mathcal S}(p-1)
       =(-1)^{\sum_{q\in\mathcal S}e_q}\lambda(p-1)B_{\mathcal S}(p-1).
\]
For fixed $Y$, the valuation event and the square exclusions at
$q\leq Y$, $q\notin\mathcal S$, are a finite union of reduced
classes modulo
\[
 \lcm\left(\prod_{q\in\mathcal S}q^{e_q+1},
           \prod_{\substack{q\leq Y\\q\notin\mathcal S}}q^2\right).
\]
Apply \eqref{eq:modecancel} with $z=-1$ and $H=1$ in each class.
Equation~\eqref{eq:squaretail} controls the remaining square factors.
Taking the limits in that order proves Corollary~\ref{cor:moree}.
The argument permits $e_q\geq2$, because squarefreeness is imposed
only outside $\mathcal S$.
